% GENERATED FILE. Edit canonical lessons, then run npm run generate:books.
% canonical-source-hash: fnv1a64:9b9f118d5d474eb0
% canonical-lessons: RU-C384-zubrit, RU-C384-zauchivat, RU-C384-progulivat, RU-C384-shpargalka, RU-C384-otvlekatsya, RU-R384-review-377-384

\chapter{Cramming and crib sheets}
\label{ch:ru-cramming-and-crib-sheets}
\hlchaptermodality{384}

\begin{chapteropening}
\textbf{By the end of this chapter:} \emph{I can talk about how students study: cramming, learning by heart, skipping classes, crib sheets and getting distracted.}

Complete the last lesson of chapter 384: I can talk about how students study: cramming, learning by heart, skipping classes, crib sheets and getting distracted.
\end{chapteropening}

\section[\texorpdfstring{\ru{зубрить} (zubrít)}{зубрить (zubrít)}]{\ru{зубрить} (zubrít) — to cram, to learn by rote}

\label{lesson:RU-C384-zubrit}



\begin{quote}
Before the new one: say the Russian for a faculty, then the Russian for a student grant.
\end{quote}

\subsection*{You'll want to know: \ru{зубрить}}
\textbf{\ru{зубрить}} (\emph{zubrít}) — \textquotedblleft{}to cram, to learn by rote\textquotedblright{}.

\begin{quote}
\ru{Перед} \ru{экзаменом} \ru{он} \ru{зубрил} \ru{всю} \ru{ночь}.
\end{quote}

\begin{quote}
\emph{Péred ekzámenom on zubríl vsyu noch.}
\end{quote}

\begin{quote}
Before the exam he crammed all night.
\end{quote}

Imperfective; present \emph{\ru{зубрю}, \ru{зубришь}}; perfective \emph{\ru{вызубрить}}. Colloquial.

\begin{cousinweb}
\textbf{\ru{зубрить}} is usually linked with \emph{\ru{зуб}} \textquotedblleft{}a tooth\textquotedblright{}, inherited from Proto-Slavic \emph{*zǫbŭ}: to gnaw away at a subject until it is learned.
\end{cousinweb}

\begin{grammarlens}[title={What you've built}]
One more word for telling, explaining and describing at B1.
\end{grammarlens}

\subsection*{Your turn}
\begin{itemize}
\raggedright
  \item \emph{Say it:} \emph{zubrít}
  \item \emph{Say it:} \emph{zubrít}, once more
  \item \emph{Say it:} a sentence of your own with \emph{zubrít}
  \item \emph{Recall:} say the Russian for a faculty, then the Russian for a student grant, then say \emph{zubrít} again
\end{itemize}

\begin{tcolorbox}[breakable,colback=teal!4,colframe=teal!35!black,title={Before you move on}]
What does \ru{зубрить} mean? (\textquotedblleft{}To cram, to learn by rote\textquotedblright{}.) Say it once more.
\end{tcolorbox}
\section[\texorpdfstring{\ru{заучивать} (zaúchivat)}{заучивать (zaúchivat)}]{\ru{заучивать} (zaúchivat) — to memorise, to learn by heart}

\label{lesson:RU-C384-zauchivat}



\begin{quote}
Before the new one: say the Russian for a student grant, then the Russian for to cram.
\end{quote}

\subsection*{You'll want to know: \ru{заучивать}}
\textbf{\ru{заучивать}} (\emph{zaúchivat}) — \textquotedblleft{}to memorise, to learn by heart\textquotedblright{}.

\begin{quote}
\ru{Мы} \ru{заучиваем} \ru{стихи} \ru{к} \ru{уроку}.
\end{quote}

\begin{quote}
\emph{My zaúchivayem stikhí k uróku.}
\end{quote}

\begin{quote}
We learn poems by heart for the lesson.
\end{quote}

Imperfective; perfective \emph{\ru{заучить}}. \emph{\ru{Заучивать} \ru{слова}} — to memorise words.

\begin{cousinweb}
\textbf{\ru{заучивать}} is built from \emph{\ru{за}-} and \emph{\ru{учить}} \textquotedblleft{}to teach, to learn\textquotedblright{}, inherited from Proto-Slavic \emph{*učiti}: to learn a thing until it sticks.
\end{cousinweb}

\begin{grammarlens}[title={What you've built}]
One more word for telling, explaining and describing at B1.
\end{grammarlens}

\subsection*{Your turn}
\begin{itemize}
\raggedright
  \item \emph{Say it:} \emph{zaúchivat}
  \item \emph{Say it:} \emph{zaúchivat}, once more
  \item \emph{Say it:} a sentence of your own with \emph{zaúchivat}
  \item \emph{Recall:} say the Russian for a student grant, then the Russian for to cram, then say \emph{zaúchivat} again
\end{itemize}

\begin{tcolorbox}[breakable,colback=teal!4,colframe=teal!35!black,title={Before you move on}]
What does \ru{заучивать} mean? (\textquotedblleft{}To memorise, to learn by heart\textquotedblright{}.) Say it once more.
\end{tcolorbox}
\section[\texorpdfstring{\ru{прогуливать} (progúlivat)}{прогуливать (progúlivat)}]{\ru{прогуливать} (progúlivat) — to skip (classes), to play truant}

\label{lesson:RU-C384-progulivat}



\begin{quote}
Before the new one: say the Russian for to cram, then the Russian for to memorise.
\end{quote}

\subsection*{You'll want to know: \ru{прогуливать}}
\textbf{\ru{прогуливать}} (\emph{progúlivat}) — \textquotedblleft{}to skip (classes), to play truant\textquotedblright{}.

\begin{quote}
\ru{Он} \ru{часто} \ru{прогуливает} \ru{лекции}.
\end{quote}

\begin{quote}
\emph{On chásto progúlivayet léktsii.}
\end{quote}

\begin{quote}
He often skips lectures.
\end{quote}

Imperfective; perfective \emph{\ru{прогулять}}. \emph{\ru{Прогул}} — a missed day.

\begin{cousinweb}
\textbf{\ru{прогуливать}} is built from \emph{\ru{про}-} \textquotedblleft{}past, through\textquotedblright{} and \emph{\ru{гулять}} \textquotedblleft{}to walk, to stroll\textquotedblright{}: to stroll past the lesson instead of going in.
\end{cousinweb}

\begin{grammarlens}[title={What you've built}]
One more word for telling, explaining and describing at B1.
\end{grammarlens}

\subsection*{Your turn}
\begin{itemize}
\raggedright
  \item \emph{Say it:} \emph{progúlivat}
  \item \emph{Say it:} \emph{progúlivat}, once more
  \item \emph{Say it:} a sentence of your own with \emph{progúlivat}
  \item \emph{Recall:} say the Russian for to cram, then the Russian for to memorise, then say \emph{progúlivat} again
\end{itemize}

\begin{tcolorbox}[breakable,colback=teal!4,colframe=teal!35!black,title={Before you move on}]
What does \ru{прогуливать} mean? (\textquotedblleft{}To skip (classes), to play truant\textquotedblright{}.) Say it once more.
\end{tcolorbox}
\section[\texorpdfstring{\ru{шпаргалка} (shpargálka)}{шпаргалка (shpargálka)}]{\ru{шпаргалка} (shpargálka) — a crib sheet}

\label{lesson:RU-C384-shpargalka}



\begin{quote}
Before the new one: say the Russian for to memorise, then the Russian for to skip.
\end{quote}

\subsection*{You'll want to know: \ru{шпаргалка}}
\textbf{\ru{шпаргалка}} (\emph{shpargálka}) — \textquotedblleft{}a crib sheet\textquotedblright{}.

\begin{quote}
\ru{На} \ru{экзамене} \ru{у} \ru{него} \ru{нашли} \ru{шпаргалку}.
\end{quote}

\begin{quote}
\emph{Na ekzámene u negó nashlí shpargálku.}
\end{quote}

\begin{quote}
They found a crib sheet on him during the exam.
\end{quote}

Feminine. Colloquial; students often say \emph{\ru{шпора}}.

\begin{cousinweb}
\textbf{\ru{шпаргалка}} is usually traced to Polish \emph{szpargał} \textquotedblleft{}an old scrap of paper\textquotedblright{}; the further origin is not settled.
\end{cousinweb}

\begin{grammarlens}[title={What you've built}]
One more word for telling, explaining and describing at B1.
\end{grammarlens}

\subsection*{Your turn}
\begin{itemize}
\raggedright
  \item \emph{Say it:} \emph{shpargálka}
  \item \emph{Say it:} \emph{shpargálka}, once more
  \item \emph{Say it:} a sentence of your own with \emph{shpargálka}
  \item \emph{Recall:} say the Russian for to memorise, then the Russian for to skip, then say \emph{shpargálka} again
\end{itemize}

\begin{tcolorbox}[breakable,colback=teal!4,colframe=teal!35!black,title={Before you move on}]
What does \ru{шпаргалка} mean? (\textquotedblleft{}A crib sheet\textquotedblright{}.) Say it once more.
\end{tcolorbox}
\section[\texorpdfstring{\ru{отвлекаться} (otvlekátsya)}{отвлекаться (otvlekátsya)}]{\ru{отвлекаться} (otvlekátsya) — to get distracted}

\label{lesson:RU-C384-otvlekatsya}



\begin{quote}
Before the new one: say the Russian for to skip, then the Russian for a crib sheet.
\end{quote}

\subsection*{You'll want to know: \ru{отвлекаться}}
\textbf{\ru{отвлекаться}} (\emph{otvlekátsya}) — \textquotedblleft{}to get distracted\textquotedblright{}.

\begin{quote}
\ru{Когда} \ru{я} \ru{работаю}, \ru{я} \ru{часто} \ru{отвлекаюсь} \ru{на} \ru{телефон}.
\end{quote}

\begin{quote}
\emph{Kogdá ya rabótayu, ya chásto otvlekáyus na telefón.}
\end{quote}

\begin{quote}
When I work, I often get distracted by my phone.
\end{quote}

Imperfective; perfective \emph{\ru{отвлечься}}. \emph{\ru{Отвлекаться} \ru{на}} + accusative — to be distracted by.

\begin{cousinweb}
\textbf{\ru{отвлекаться}} is built from \emph{\ru{от}-} \textquotedblleft{}away\textquotedblright{} and \emph{\ru{влечь}} \textquotedblleft{}to drag, to pull\textquotedblright{}, an inherited Slavic verb: to be pulled away from your work.
\end{cousinweb}

\begin{grammarlens}[title={What you've built}]
One more word for telling, explaining and describing at B1.
\end{grammarlens}

\subsection*{Your turn}
\begin{itemize}
\raggedright
  \item \emph{Say it:} \emph{otvlekátsya}
  \item \emph{Say it:} \emph{otvlekátsya}, once more
  \item \emph{Say it:} a sentence of your own with \emph{otvlekátsya}
  \item \emph{Recall:} say the Russian for to skip, then the Russian for a crib sheet, then say \emph{otvlekátsya} again
\end{itemize}

\begin{tcolorbox}[breakable,colback=teal!4,colframe=teal!35!black,title={Before you move on}]
What does \ru{отвлекаться} mean? (\textquotedblleft{}To get distracted\textquotedblright{}.) Say it once more.
\end{tcolorbox}
\section[(dialogue)]{Words from chapters 377–384 — a review}

\label{lesson:RU-R384-review-377-384}



\begin{quote}
Say the Russian for to postpone and to change one’s mind.
\end{quote}

\subsection*{Your turn}
Say these aloud:

\begin{itemize}
\raggedright
  \item to postpone, to change one’s mind, to refuse, to set, to be going to
  \item to arrange, an arrangement, an option, an alternative, final
  \item spare, in any case, to hesitate, to make up one’s mind, depending on
  \item to predict, to guess, a deadline, permanent, temporary
  \item an exam period, a pass-or-fail test, a seminar, lecture notes, a point
  \item to fail, to retake, a resit, preparation, to prepare oneself
  \item admission, a university applicant, entrance, a faculty, a student grant
  \item to cram, to memorise, to skip, a crib sheet, to get distracted
\end{itemize}

\begin{tcolorbox}[breakable,colback=teal!4,colframe=teal!35!black,title={Before you move on}]
To postpone? (\textbf{\ru{переносить}}.) To change one’s mind? (\textbf{\ru{передумывать}}.) To refuse? (\textbf{\ru{отказываться}}.) To set? (\textbf{\ru{назначать}}.) To be going to? (\textbf{\ru{собираться}}.) To arrange? (\textbf{\ru{договариваться}}.) An arrangement? (\textbf{\ru{договорённость}}.) An option? (\textbf{\ru{вариант}}.) An alternative? (\textbf{\ru{альтернатива}}.) Final? (\textbf{\ru{окончательный}}.) Spare? (\textbf{\ru{запасной}}.) In any case? (\textbf{\ru{в} \ru{любом} \ru{случае}}.) To hesitate? (\textbf{\ru{колебаться}}.) To make up one’s mind? (\textbf{\ru{решаться}}.) Depending on? (\textbf{\ru{в} \ru{зависимости} \ru{от}}.) To predict? (\textbf{\ru{предсказывать}}.) To guess? (\textbf{\ru{угадывать}}.) A deadline? (\textbf{\ru{дедлайн}}.) Permanent? (\textbf{\ru{постоянный}}.) Temporary? (\textbf{\ru{временный}}.) An exam period? (\textbf{\ru{сессия}}.) A pass-or-fail test? (\textbf{\ru{зачёт}}.) A seminar? (\textbf{\ru{семинар}}.) Lecture notes? (\textbf{\ru{конспект}}.) A point? (\textbf{\ru{балл}}.) To fail? (\textbf{\ru{провалить}}.) To retake? (\textbf{\ru{пересдавать}}.) A resit? (\textbf{\ru{пересдача}}.) Preparation? (\textbf{\ru{подготовка}}.) To prepare oneself? (\textbf{\ru{готовиться}}.) Admission? (\textbf{\ru{поступление}}.) A university applicant? (\textbf{\ru{абитуриент}}.) Entrance? (\textbf{\ru{вступительный}}.) A faculty? (\textbf{\ru{факультет}}.) A student grant? (\textbf{\ru{стипендия}}.) To cram? (\textbf{\ru{зубрить}}.) To memorise? (\textbf{\ru{заучивать}}.) To skip? (\textbf{\ru{прогуливать}}.) A crib sheet? (\textbf{\ru{шпаргалка}}.) To get distracted? (\textbf{\ru{отвлекаться}}.)
\end{tcolorbox}
